\chapter{Shell Programming, Automation Scripting \& Control Flow}
\label{chap:shell-programming}

\section{Anatomy of an Industrial Shell Script}
A shell script is an executable ASCII text file housing a sequence of commands executed by the shell interpreter. Shell scripting transforms manual terminal commands into repeatable, deterministic, automated systems.

\subsection{The Shebang (\texttt{\#!/bin/bash})}
The very first line of a script must begin with the two-byte magic number \texttt{\#!} (shebang):
\begin{scriptbox}[Basic Shell Script Skeleton]
#!/bin/bash
# Description: Automated system health check script
# Author: DevOps Engineering Team

echo "System Health Inspection: $(date)"
\end{scriptbox}
\begin{itemize}[leftmargin=1.5em]
    \item When an executable script is loaded via the \texttt{execve()} system call, the Linux kernel inspects the first two bytes (\texttt{0x23 0x21}). If present, the kernel extracts the absolute path to the interpreter (\texttt{/bin/bash}) and executes it, passing the script path as an argument.
    \item \textbf{Portability Pattern:} Using \texttt{\#!/usr/bin/env bash} resolves the interpreter path dynamically from the user's \texttt{\$PATH}.
\end{itemize}

\subsection{Script Execution Environments}
\begin{table}[htbp]
\centering
\small
\begin{tabularx}{\linewidth}{l X X}
\toprule
\textbf{Invocation Method} & \textbf{Process Execution Model} & \textbf{Impact on Environment} \\
\midrule
\texttt{./script.sh} & Kernel spawns a child subshell. Requires read (\texttt{r}) and execute (\texttt{x}) permissions. & Variables and working directory changes expire upon script exit. \\
\texttt{bash script.sh} & Shell explicitly executes binary. Requires read (\texttt{r}) permission only. & Isolated child subshell; parent environment unchanged. \\
\texttt{source script.sh} or \texttt{. script.sh} & Executed directly \textbf{within the current shell process}. No child process is created. & Changes to variables, aliases, functions, and directory paths \textbf{persist} in the active shell. \\
\bottomrule
\end{tabularx}
\caption{Comparison of Shell Script Execution Invocations.}
\label{tab:script-exec}
\end{table}

\section{Variables, Quoting Rules \& Parameter Expansion}
Variables in Bash are untyped character strings by default.

\subsection{Assignment Syntax and Quoting}
\begin{warningbox}[Variable Assignment Whitespace Rule]
Variable assignment \textbf{must not contain spaces} around the equals sign:
\begin{itemize}
    \item \texttt{NAME="Linux"} $\to$ Valid variable assignment.
    \item \texttt{NAME = "Linux"} $\to$ Shell interprets \texttt{NAME} as a command and \texttt{=} as an argument, causing a \texttt{command not found} runtime failure.
\end{itemize}
\end{warningbox}

\begin{itemize}[leftmargin=1.5em]
    \item \textbf{Single Quotes (\texttt{'...'}):} Preserve the strict literal value of every character within. No variable expansion, escape sequence, or command substitution is performed.
    \item \textbf{Double Quotes (\texttt{"..."}):} Preserve literal values of most characters while allowing parameter expansion (\texttt{\$VAR}), command substitution (\texttt{\$(...)}), and arithmetic evaluation (\texttt{\$(...)}).
    \item \textbf{Command Substitution:} The syntax \texttt{\$(command)} captures the standard output stream of an external command into a variable (e.g., \texttt{UPTIME=\$(uptime -p)}).
\end{itemize}

\subsection{Positional Parameters \& Special Built-in Variables}
\begin{table}[htbp]
\centering
\small
\begin{tabularx}{\linewidth}{l X}
\toprule
\textbf{Variable} & \textbf{Semantics and Purpose} \\
\midrule
\texttt{\$0} & Name of the script currently executing (as invoked). \\
\texttt{\$1 ... \$9} & First through ninth command-line positional arguments. \\
\texttt{\$\{10\}} & Tenth positional argument (requires curly braces for two or more digits). \\
\texttt{\$\#} & Total count of positional arguments passed to the script. \\
\texttt{\$@} & All positional arguments as separate quoted words (\texttt{"\$1" "\$2" ...}). \\
\texttt{\$*} & All positional arguments combined into a single word separated by \texttt{IFS}. \\
\texttt{\$?} & Exit status code of the most recently executed foreground command. \\
\texttt{\$\$} & Process ID (PID) of the currently running script instance. \\
\texttt{\$!} & Process ID of the most recently spawned background job. \\
\bottomrule
\end{tabularx}
\caption{POSIX and Bash Special Built-in Variables.}
\label{tab:special-vars}
\end{table}

\section{Operators \& Conditional Testing: [ ] vs [[ ]]}
Conditionals evaluate exit statuses. The \texttt{test} command and its alias \texttt{[ ... ]} are POSIX standards; \texttt{[[ ... ]]} is an enhanced Bash keyword that prevents word splitting and pathname expansion bugs.

\begin{table}[htbp]
\centering
\small
\begin{tabularx}{\linewidth}{l X l X}
\toprule
\textbf{Numeric Test} & \textbf{Condition Checked} & \textbf{String Test} & \textbf{Condition Checked} \\
\midrule
\texttt{-eq} & Equal to ($a = b$) & \texttt{=} or \texttt{==} & Strings are identical \\
\texttt{-ne} & Not equal to ($a \neq b$) & \texttt{!=} & Strings are distinct \\
\texttt{-lt} & Less than ($a < b$) & \texttt{-z str} & String length is zero (empty) \\
\texttt{-le} & Less than or equal ($a \le b$) & \texttt{-n str} & String length is non-zero \\
\texttt{-gt} & Greater than ($a > b$) & \texttt{<}, \texttt{>} & Lexicographic ASCII order \\
\texttt{-ge} & Greater than or equal ($a \ge b$) & & \\
\bottomrule
\end{tabularx}
\caption{Relational Operators for Numeric and String Comparisons.}
\label{tab:rel-ops}
\end{table}

\begin{table}[htbp]
\centering
\small
\begin{tabularx}{\linewidth}{l X l X}
\toprule
\textbf{File Operator} & \textbf{Condition Evaluated} & \textbf{File Operator} & \textbf{Condition Evaluated} \\
\midrule
\texttt{-e file} & Path exists (file, directory, socket) & \texttt{-r file} & File exists and is readable \\
\texttt{-f file} & Path exists and is a \textbf{regular file} & \texttt{-w file} & File exists and is writable \\
\texttt{-d dir} & Path exists and is a \textbf{directory} & \texttt{-x file} & File exists and is executable \\
\texttt{-s file} & File exists and size is $> 0$ bytes & \texttt{-L file} & Path is a symbolic link \\
\bottomrule
\end{tabularx}
\caption{Filesystem Attribute Test Operators.}
\label{tab:file-ops}
\end{table}

\section{Control Structures \& Decision-Making Constructs}
\subsection{The if-elif-else-fi Branching Construct}
\begin{scriptbox}[Decision Branching Example]
#!/bin/bash
TARGET="/var/log/syslog"

if [ ! -e "$TARGET" ]; then
    echo "Error: File $TARGET does not exist." >&2
    exit 1
elif [ -r "$TARGET" ] && [ -s "$TARGET" ]; then
    echo "Target is present, readable, and non-empty."
    LINE_COUNT=$(wc -l < "$TARGET")
    echo "Total log lines: $LINE_COUNT"
else
    echo "Warning: Target exists but is unreadable or empty."
fi
\end{scriptbox}

\subsection{The case-esac Pattern Matching Construct}
The \texttt{case} statement provides structured multi-way branching against globbing patterns:
\begin{scriptbox}[Case Statement Example]
#!/bin/bash
case "$1" in
    start)
        echo "Starting application service daemon..."
        ;;
    stop)
        echo "Stopping application service daemon..."
        ;;
    restart|reload)
        echo "Restarting service..."
        ;;
    *)
        echo "Usage: $0 {start|stop|restart}" >&2
        exit 2
        ;;
esac
\end{scriptbox}

\section{Looping Constructs: for, while, until}
\subsection{The for Loop}
\begin{itemize}[leftmargin=1.5em]
    \item \textbf{Iterating over a collection:}
\begin{scriptbox}[Iterating Collections]
for file in /etc/*.conf; do
    echo "Configuring file: $file"
done
\end{scriptbox}
    \item \textbf{C-Style Three-Expression Loop:}
\begin{scriptbox}[C-Style Loop]
for (( i=1; i<=10; i++ )); do
    printf "Square of \%d is \%d\n" "$i" "$(( i * i ))"
done
\end{scriptbox}
\end{itemize}

\subsection{The while Loop \& Streaming File Processing}
The standard pattern for reading files line-by-line safely:
\begin{scriptbox}[Line-by-Line File Processing]
#!/bin/bash
INPUT_FILE="/etc/hosts"

while IFS= read -r line || [ -n "$line" ]; do
    # Skip comment lines
    [[ "$line" =~ ^# ]] && continue
    echo "Processing host entry: $line"
done < "$INPUT_FILE"
\end{scriptbox}

\section{Modular Programming: Functions and Arrays}
\subsection{Bash Functions and Variable Scoping}
Functions allow isolating logic and parameterizing scripts:
\begin{scriptbox}[Modular Function Declaration]
#!/bin/bash

# Function definition
backup_directory() {
    local SRC="$1"
    local DEST="$2"

    if [ ! -d "$SRC" ]; then
        echo "Source directory $SRC not found!" >&2
        return 1
    fi

    tar -czf "$DEST/backup_$(date +\%Y\%m\%d).tar.gz" -C "$SRC" .
    return 0
}

# Function invocation
backup_directory "/var/www/html" "/backup"
if [ $? -eq 0 ]; then
    echo "Backup completed successfully."
fi
\end{scriptbox}

\subsection{Bash Arrays: Indexed and Associative}
\begin{itemize}[leftmargin=1.5em]
    \item \textbf{Indexed Arrays:}
\begin{scriptbox}[Indexed Arrays]
servers=("alpha" "beta" "gamma")
echo "First server: ${servers[0]}"
echo "Total servers: ${#servers[@]}"
servers+=("delta")   # Append element
\end{scriptbox}
    \item \textbf{Associative Arrays (Hash Maps):}
\begin{scriptbox}[Associative Arrays]
declare -A ports
ports["http"]=80
ports["https"]=443
ports["ssh"]=22
echo "SSH Port is ${ports["ssh"]}"
\end{scriptbox}
\end{itemize}

\section{Defensive Scripting Standards}
Production scripts should enforce strict error handling at the entry point:
\begin{scriptbox}[Defensive Entry Template]
#!/bin/bash
set -euo pipefail
# -e: Exit immediately if any command exits with a non-zero status.
# -u: Treat unset variables as an error and exit immediately.
# -o pipefail: Pipeline fails if any command inside encounters an error.

trap 'rm -rf "$TMP_DIR"; echo "Cleanup complete."' EXIT
TMP_DIR=$(mktemp -d)
\end{scriptbox}

\section{Chapter Review \& Analytical Questions}
\begin{enumerate}[leftmargin=1.5em]
    \item Explain the mechanics of the shebang line (\texttt{\#!/bin/bash}) and describe what occurs when it is omitted.
    \item Why must variables be enclosed in double quotes when testing string equality (e.g., \texttt{[ "\$VAR" = "prod" ]})?
    \item Distinguish between the special variables \texttt{\$@} and \texttt{\$*}.
    \item Write a complete Bash script that accepts a directory path as argument, verifies it exists, and counts the number of regular files vs subdirectories inside.
    \item Explain why \texttt{set -euo pipefail} is designated as the gold standard for reliable DevOps shell scripting.
\end{enumerate}
